\chapter{Diagnostic Transition and Integrated Architecture}
\label{chap:architecture_retained_method}

% Writing objective:
% This chapter should bridge the retained controller line and the newer integrated solver line.
% It should first explain the large-instance diagnostic finding, then introduce the integrated
% controller-routing-repair architecture that follows from that diagnosis.

% ============================================================
\section{Retained execution architecture in the controller-centered line}
% Write here:
% - Briefly summarize the operational architecture behind 22.03controller_line.
% - Explain controller + daily backend + experiment workflow at a high level.
% - Keep this concise; the focus is on why this architecture was useful and what it exposed.
%
% Good source material:
% - 22.03controller_line/docs/architecture.md
% - 22.03controller_line/scripts/
% - 22.03controller_line/code/

\section{Diagnostic audit on large-scale instances}
% Write here:
% - Introduce West as the critical audit case.
% - Explain why a deeper diagnosis became necessary.
% - Present the logic of bucket-level diagnostics.
%
% Good source material:
% - 22.04fresh_solver/docs/OR支线实验_v7.md
% - 22.04fresh_solver/results/diagnostics/west_or_v7_picking_buckets.csv
% - 22.04fresh_solver/results/diagnostics/west_or_v7_picking_summary.json
%
% Key message:
% The bottleneck is not all-day capacity shortage, but a sharply localized morning picking wave.

\section{Bottleneck migration: from admission pressure to depot-wave pressure}
% Write here:
% - State the key scientific finding cleanly:
%   the dominant overload in West is concentrated around 05:30--06:15,
%   and almost all overload is before noon.
% - Explain why this changes the problem interpretation.
%
% Good source material:
% - 22.04fresh_solver/results/diagnostics/west_or_v7_picking_summary.json
% - 22.04fresh_solver/docs/OR支线_V8动态shadow_price阶段总结_v1.md
% - Desktop summary markdown on the Desktop

\section{Why controller-only smoothing and generic backends are insufficient}
% Write here:
% - Explain why high-level admission control alone cannot fully smooth a localized bucket peak.
% - Explain why generic routing backends do not naturally internalize bucket-level warehouse pressure.
% - This is the motivation section for the integrated architecture.
%
% Good source material:
% - 22.04fresh_solver/docs/最优算法架构.md
% - 22.04fresh_solver/docs/22.03controller_line_vs_22.04fresh_solver_关系与口径说明.md
%
% Important:
% Keep the tone analytical, not dismissive of 22.03.

\section{Integrated controller-routing-repair architecture}
% Write here:
% - Present the three-layer architecture explicitly:
%   controller + routing + depot-repair.
% - Explain the role of each layer and how information flows between them.
%
% Good source material:
% - 22.04fresh_solver/docs/最优算法架构.md
% - 22.04fresh_solver/README.md
% - 22.04fresh_solver/src/algorithms/fresh_solver/julia/README.md
%
% Suggested subpoints to describe:
% - controller: admission / protection / clipping
% - routing: route construction under dynamic costs
% - repair: bucket/depot pressure mitigation

\section{Implementation shift: from Python-centered backend to Julia-first solver}
% Write here:
% - Explain why a Julia-first implementation became useful for the integrated line.
% - Focus on computational and architectural reasons.
% - Do not oversell this as a language war.
%
% Good source material:
% - 22.04fresh_solver/src/algorithms/fresh_solver/julia/README.md
% - 22.04fresh_solver/src/algorithms/fresh_solver/julia/src/FreshSolver.jl
% - 22.04fresh_solver/src/algorithms/fresh_solver/julia/src/io.jl
